\documentclass[10pt,a4paper]{amsart}
\usepackage[margin=19mm]{geometry}
\usepackage{amsmath,amssymb,mathtools}
\usepackage[scr=rsfs,cal=euler]{mathalfa}
\usepackage{xfrac}
\usepackage[unicode,hidelinks,bookmarks=false]{hyperref}
\newcommand{\ie}{i.e.\ }
\newcommand{\cf}{cf.\ }
\newcommand{\resp}{respectively\ }
\newcommand{\Subsection}{\subsection}
\providecommand{\nobreakdash}{\nobreak\hbox{-}}
\setlength{\emergencystretch}{2em}
\multlinegap0pt
\usepackage{bm}
\usepackage{bbm}
\usepackage{dsfont}
\usepackage{xspace}
\usepackage{cancel}
\usepackage{mathtools}
\usepackage{amsthm}
\makeatletter
\@ifundefined{definition}{\newtheorem{definition}{Definition}}{}
\@ifundefined{lemma}{\newtheorem{lemma}[definition]{Lemma}}{}
\makeatother
\def\R{{\mathbb R}}
\def\T{{\mathbb T}}
\def\Z{{\mathbb Z}}
\newcommand{\ft}{{\mathcal{F}}}
\newcommand{\lt}{{L^2(\R)}}
\title[Compactly supported Gabor orthonormal bases]{Compactly supported Gabor orthonormal bases}
\author{Lukas Liehr}
\date{Source: arXiv:2605.29984v1 (CC BY 4.0)}
\begin{document}
\maketitle

\noindent\emph{Source and licence.} Compactly supported Gabor orthonormal bases. Version arXiv:2605.29984v1 (CC BY 4.0), distributed under the Creative Commons Attribution 4.0 International licence, as recorded in the arXiv licence metadata for this exact version and shown on the abstract page https://arxiv.org/abs/2605.29984v1. Full rights notice: licenses/arxiv-2605.29984-CC-BY-4.0.txt.\par\smallskip

\noindent\textit{Editorial scope.} Counted unit: the Lemma whose source label is \texttt{lem:direct-fibre-form} in the article (source0.tex lines 133--144, source label \texttt{lem:direct-fibre-form}), delivered together with the article's own complete original proof of it (source0.tex lines 146--234, one \texttt{proof} environment of 89 lines). No mathematics is written, repaired or re-derived: statement and proof are the article's own text line for line. Required context retained uncounted: none. Every definition, display and label of those lines is retained unchanged. Omitted from this package and not redistributed: 1--132, 145--145, 235--843. Nothing inside a retained excerpt is omitted or altered: every retained body is delivered exactly as the article's lines are written, so no omission receipt of any kind accompanies this package. Citations: the retained excerpts carry no citation command at all, so no bibliography entry of the article is transcribed and no reference is redistributed. Reference adaptation: none was needed and none was made; every reference inside the delivered unit resolves inside it, so no unresolved reference or citation is produced. Printed slips kept literal: no printed word, symbol, hypothesis or number of the counted statement or of its proof was altered. Layout and class: the article's own class is \texttt{amsart} with packages including inputenc, amsmath, amsthm, amssymb, amscd, accents, bm, amsfonts, mathtools, mathrsfs, dsfont, datetime, hyperref, colonequals; the wrapper is the lane's pinned \texttt{amsart} editorial wrapper at 10pt on A4, which restates the theorem-like environments the retained text needs and adds the article's own macro definitions that the retained text needs (\texttt{\textbackslash R}, \texttt{\textbackslash T}, \texttt{\textbackslash Z}, \texttt{\textbackslash ft}, \texttt{\textbackslash lt}), with extra wrapper packages (bm, bbm, dsfont, xspace, cancel, mathtools). The delivered numbering is the unit's own. Assets: no retained span contains a figure environment, a graphics-inclusion command, a PDF-page inclusion, a table or a TikZ picture; no image file is copied into this package, the package declares no asset dependency and no image file is redistributed with it. Licence: version arXiv:2605.29984v1 of this article is distributed under the Creative Commons Attribution 4.0 International licence, as recorded in the arXiv licence metadata for this exact version and shown on the abstract page https://arxiv.org/abs/2605.29984v1. A full rights notice is delivered at licenses/arxiv-2605.29984-CC-BY-4.0.txt. Characters: every retained excerpt is pure ASCII, so the delivered engine, XeLaTeX with its default Latin Modern text font, reports no missing character.
\begin{lemma}\label{lem:direct-fibre-form}
Let $g \in \lt$ and let $\alpha \in \R$ such that $\left\{M_nT_{m+\alpha n}g:m,n\in\mathbb Z\right\}$ is an orthonormal system in $L^2(\R)$. Define a sequence $\{ \gamma_k \}_{k \in \Z}$ by $\gamma_k=e^{\pi i\alpha k(k-1)}$. Moreover, define for a.e. $\omega\in\mathbb R$,
$$
D(\omega,\theta) = \sum_{k\in\mathbb Z} \gamma_k\widehat g(k+\omega)e^{2\pi i k\theta}
$$
where the series is interpreted as an $L^2(\T)$-Fourier series in $\theta$. Then 
\begin{equation}\label{eq:direct-fibre-covariance}
D(\omega+1,\theta) = e^{2\pi i(\alpha-\theta)}D(\omega,\theta-\alpha)
\end{equation}
for a.e. $(\omega,\theta) \in \R \times \T$.
Moreover, $|D(\omega,\theta)|=1$ for a.e. $(\omega,\theta)\in\mathbb R\times\T$.
\end{lemma}

\begin{proof}
Define the sequences $h(\omega)=\{h_k(\omega)\}_{k\in\Z}$ and $b(\omega)=\{b_k(\omega)\}_{k\in\Z}$ by
$$
h_k(\omega)=\widehat g(k+\omega), \quad b_k(\omega)=\gamma_k h_k(\omega).
$$
Since $\widehat g\in L^2(\mathbb R)$, it follows that for a.e. $\omega \in [0,1]$, the sequence $h(\omega)$ belongs to $\ell^2(\Z)$. Since $|\gamma_k|=1$, the same is true for $b(\omega)$.
For $n\in\mathbb Z$, define
$$
r_n(\omega) = \sum_{k\in\mathbb Z} b_{k-n}(\omega)\overline{b_k(\omega)}.
$$
This series is absolutely convergent for a.e. $\omega \in [0,1]$. Indeed, by the Cauchy-Schwarz inequality
$$
\sum_{k\in\mathbb Z} |b_{k-n}(\omega)||b_k(\omega)| \leq \left(\sum_{k\in\mathbb Z}|b_{k-n}(\omega)|^2\right)^{1/2} \left(\sum_{k\in\mathbb Z}|b_k(\omega)|^2\right)^{1/2} = \sum_{k\in\mathbb Z}|b_k(\omega)|^2.
$$
Since $|b_k(\omega)|=|h_k(\omega)|$, it follows from the estimate
$$
\int_0^1 |r_n(\omega)| \, d\omega \leq \int_0^1  \sum_{k\in\mathbb Z}|h_k(\omega)|^2 \, d\omega = \| g \|^2
$$
that $r_n \in L^1(0,1)$. The identity
$$
\ft M_nT_{m+\alpha n}g(\xi) = e^{-2\pi i(m+\alpha n)(\xi-n)} \ft g(\xi-n),
$$
in combination with orthonormality and Plancherel's theorem gives
$$
\delta_{m,0}\delta_{n,0} = \int_{\mathbb R} e^{-2\pi i(m+\alpha n)(\xi-n)} \widehat g(\xi-n)\overline{\widehat g(\xi)}\,d\xi .
$$
Equivalently, we have
$$
\delta_{m,0}\delta_{n,0} = \int_{\mathbb R} e^{-2\pi i(m+\alpha n)\xi} \widehat g(\xi-n)\overline{\widehat g(\xi)}\,d\xi .
$$
Periodization of the integral and using the definition of
$h(\omega)$, we get
$$
\delta_{m,0}\delta_{n,0} = \int_0^1 e^{-2\pi i m\omega} e^{-2\pi i\alpha n\omega} e^{2\pi i\alpha n^2} \sum_{k\in\mathbb Z} e^{-2\pi i\alpha nk} h_{k-n}(\omega)\overline{h_k(\omega)} \,d\omega .
$$
The elementary identity $\gamma_{k-n}\overline{\gamma_k} = e^{-2\pi i\alpha nk}e^{\pi i\alpha n(n+1)}$ shows that the series under the previous integral satisfies the relation
$$
\sum_{k\in\mathbb Z} e^{-2\pi i\alpha nk} h_{k-n}(\omega)\overline{h_k(\omega)} = e^{-\pi i\alpha n(n+1)}r_n(\omega).
$$
Consequently,
$$
\delta_{m,0}\delta_{n,0} = e^{\pi i\alpha n(n-1)} \int_0^1 e^{-2\pi i m\omega} e^{-2\pi i\alpha n\omega} r_n(\omega)\,d\omega.
$$
Since the prefactor is nonzero, we obtain
$$
\int_0^1 e^{-2\pi i m\omega} \bigl(e^{-2\pi i\alpha n\omega}r_n(\omega)\bigr) \,d\omega = \delta_{m,0}\delta_{n,0}, \quad  m\in\mathbb Z.
$$
We have shown above that $r_n \in L^1(0,1)$. Therefore the function
$$
e^{-2\pi i\alpha n\omega}r_n(\omega)
$$
belongs to $L^1(0,1)$. Hence, by uniqueness of Fourier coefficients,
$$
e^{-2\pi i\alpha n\omega}r_n(\omega)=\delta_{n,0}
$$
for a.e. $\omega\in[0,1]$. Thus, after removing one null set independent of $n$,
$$
r_n(\omega)=\delta_{n,0}, \quad n \in \mathbb Z,
$$
for a.e. $\omega\in[0,1]$.
For such an $\omega$, define
$$
D(\omega,\theta) = \sum_{k\in\mathbb Z}b_k(\omega)e^{2\pi i k\theta}
$$
as an $L^2(\T)$-Fourier series. Then $|D(\omega,\cdot)|^2\in L^1(\T)$. Its Fourier coefficients are the autocorrelations of $b(\omega)$,
$$
\int_0^1 |D(\omega,\theta)|^2e^{2\pi i n\theta}\,d\theta = \sum_{k\in\mathbb Z} b_{k-n}(\omega)\overline{b_k(\omega)} = r_n(\omega) = \delta_{n,0}.
$$
Thus all Fourier coefficients of the $L^1$-function $|D(\omega,\cdot)|^2$ agree with those of the constant function $1$. By uniqueness of Fourier coefficients, $|D(\omega,\theta)|^2=1$ for a.e. $\theta\in \T $. Consequently, $|D(\omega,\theta)|=1$ for a.e. $(\omega,\theta)\in[0,1]\times \T$.
Further, we have
$$
\begin{aligned}
D(\omega+1,\theta)
&=
\sum_{k\in\mathbb Z}
\gamma_k\widehat g(k+\omega+1)e^{2\pi i k\theta}  \\
&=
\sum_{j\in\mathbb Z}
\gamma_{j-1}\widehat g(j+\omega)e^{2\pi i(j-1)\theta}  \\
&=
e^{2\pi i(\alpha-\theta)}
\sum_{j\in\mathbb Z}
\gamma_j\widehat g(j+\omega)e^{2\pi i j(\theta-\alpha)}  \\
&=
e^{2\pi i(\alpha-\theta)}D(\omega,\theta-\alpha).
\end{aligned}
$$
Finally, since the factor $e^{2\pi i(\alpha-\theta)}$ has modulus one, this identity transfers $|D|=1$ from $[0,1]\times \T$ to every shift $(r+[0,1])\times \T$ with $r\in\mathbb Z$, implying that $|D(\omega,\theta)|=1$ for a.e. $(\omega,\theta)\in\mathbb R\times \T$.
\end{proof}

\end{document}
